\begin{theorem}[Chebyshev’s Sum Inequality]
Let  
\[
a_{1}\le a_{2}\le\cdots\le a_{n},\qquad   
b_{1}\le b_{2}\le\cdots\le b_{n}
\]
be real numbers. Then  
\[
\boxed{\,n\sum_{i=1}^{n}a_i b_i\;\ge\;
\Bigl(\sum_{i=1}^{n}a_i\Bigr)\Bigl(\sum_{i=1}^{n}b_i\Bigr)\,}.
\]
\end{theorem}

\begin{proof}
\begin{enumerate}
\item \textbf{Degenerate cases.}
\begin{itemize}
\item $n=0$: the statement is vacuously true.
\item $n=1$: the inequality reads $a_{1}b_{1}\ge a_{1}b_{1}$, true with equality.
\item If the $a$–sequence (or the $b$–sequence) is constant, each factor $(a_i-a_j)$ (resp.\ $(b_i-b_j)$) is $0$, so the two sides are equal.  These are the trivial equality cases.
\end{itemize}
Hence we may assume $n\ge2$ and that at least one of the sequences is not constant.

\item \textbf{A non‑negative double sum.}  
Define  
\[
S:=\sum_{1\le i<j\le n}(a_i-a_j)(b_i-b_j).
\]
Because the sequences are non‑decreasing, for $i<j$ we have $a_i-a_j\le0$ and $b_i-b_j\le0$, whence  
\[
(a_i-a_j)(b_i-b_j)\ge0 .
\]
Summing over all $\binom{n}{2}$ pairs gives  
\[
S\ge0. \tag{1}
\]

\item \textbf{Expansion of $S$.}  Write the sum symmetrically:
\[
S=\frac12\sum_{i=1}^{n}\sum_{j=1}^{n}(a_i-a_j)(b_i-b_j).
\]
Expanding the product inside the double sum,
\begin{align*}
S&=\frac12\Bigl(
\sum_i\sum_j a_i b_i
-\sum_i\sum_j a_i b_j
-\sum_i\sum_j a_j b_i
+\sum_i\sum_j a_j b_j\Bigr)\\[2mm]
&=\frac12\Bigl(
n\sum_{i=1}^{n}a_i b_i
-\bigl(\sum_{i=1}^{n}a_i\bigr)\bigl(\sum_{j=1}^{n}b_j\bigr)
-\bigl(\sum_{j=1}^{n}a_j\bigr)\bigl(\sum_{i=1}^{n}b_i\bigr)
+n\sum_{i=1}^{n}a_i b_i\Bigr)\\[2mm]
&=n\sum_{i=1}^{n}a_i b_i-
\Bigl(\sum_{i=1}^{n}a_i\Bigr)\Bigl(\sum_{i=1}^{n}b_i\Bigr).
\end{align*}
Thus we have the identity
\[
\sum_{1\le i<j\le n}(a_i-a_j)(b_i-b_j)
=
n\sum_{i=1}^{n}a_i b_i-
\Bigl(\sum_{i=1}^{n}a_i\Bigr)\Bigl(\sum_{i=1}^{n}b_i\Bigr). \tag{2}
\]

\item \textbf{Deriving the inequality.}  
From (1) and (2) we obtain
\[
n\sum_{i=1}^{n}a_i b_i-
\Bigl(\sum_{i=1}^{n}a_i\Bigr)\Bigl(\sum_{i=1}^{n}b_i\Bigr)\ge0,
\]
which is exactly the desired inequality.

\item \textbf{Equality condition.}  
Equality holds iff $S=0$.  Since each term $(a_i-a_j)(b_i-b_j)$ is non‑negative, $S=0$ precisely when every term vanishes, i.e.
\[
(a_i-a_j)(b_i-b_j)=0\quad\text{for all }1\le i<j\le n.
\]
Hence for each pair $(i,j)$ we must have either $a_i=a_j$ or $b_i=b_j$.  Consequently the equality cases are:

\begin{enumerate}
\item All $a_i$ are equal (the $a$‑sequence is constant);
\item All $b_i$ are equal (the $b$‑sequence is constant);
\item More generally, the index set $\{1,\dots,n\}$ can be partitioned into blocks such that on each block the $a$’s are equal and, on the same block, the $b$’s are equal.
\end{enumerate}
When both sequences are strictly increasing, only the first two possibilities can occur.

\end{enumerate}
\end{proof}